\chapter{Software Architecture Design}
\label{chap:software-architecture-design}

\section{System Architecture Overview}
\label{section:architecture-overview}
% <TIP: Provide a high-level description of your system's structure.
% Include either a Domain Model (to describe core business concepts
% and their relationships)
% or a Design Class Diagram (to describe how your code is organized
% into classes and packages),
% whichever best communicates your design.
% Accompany your diagram with a short paragraph explaining the key
% components and how they interact. />

\subsection{Domain Model}
\label{section:domain-model}
\begin{figure}[htbp]
    \centering
    \includegraphics[width=0.95\textwidth]{images/domain_model.png}
    \caption{ChatreeDeka Domain Model}
\end{figure}

ChatreeDeka's domain model is built around the CaseSimulation, which controls the overall flow of the legal dispute simulation. A User starts the process by entering the case facts, evidence, and legal claims. The system then coordinates several Agents to play different courtroom roles, such as the plaintiff, defendant, auditor, and judge. These agents interact by creating and reviewing Arguments, which can be supported by LegalProvisions and Precedents. During the simulation, the auditor checks the reliability of the arguments and records the results in AuditLogs, while the judge reviews the discussion and gives a Verdict. After the simulation is completed, the system generates a PostTrialAnalysis to explain the result, provide feedback, and suggest possible next steps.
\subsection{Architecture Design}
\label{section:architecture-design}
The software architecture follows a layered design, separating the presentation layer (user inputs and case visualization), the application layer (case orchestration and agent coordination), the AI/agent layer (role-specific reasoning and argument generation), and the data/knowledge layer (legal provisions, precedents, and audit logs). This separation keeps responsibilities clear, supports independent evolution of the AI components, and reduces coupling between the user interface and the underlying legal reasoning pipeline.

\begin{figure}[htbp]
    \centering
    \includegraphics[width=0.95\textwidth]{images/software_arch.png}
    \caption{Layered Software Architecture}
\end{figure}

\section{Software Design}
\label{section:software-design}
% <TIP: Describe the runtime behavior of your system.
% Include a Sequence Diagram for your most important use cases.
% Focus on scenarios with non-trivial interactions between components.
% If your project involves a novel algorithm,
% describe it here and explain why you designed it this way. />
\subsection{System Overview}

\textbf{ChatreeDeka} is an AI-driven platform for simulating trials. It features \textbf{4 core agents} trained via Multi-Agent Reinforcement Learning (MARL), with \textbf{hybrid user control}:

\begin{itemize}[leftmargin=*]
\item \textbf{Judge Agent}: Manual user input (rulings, objections, procedure).
\item \textbf{Defender Lawyer Agent}: Auto-generates defenses, cross-exams (MARL-trained for realism).
\item \textbf{Prosecutor Lawyer Agent}: Manual user input (case-building, exams).
\item \textbf{Case Predictor Agent}: Analyzes full sim outputs (texts, evidence, results) to predict winning probability (\textbf{\%} for each side).
\end{itemize}

\subsection{Core Workflow}

\begin{enumerate}
\item Users input case details/evidence $\rightarrow$ simulation launches.
\item \textbf{Trial}: User manually provides Judge/Prosecutor lines; Defender auto-responds in real-time mock trial.
\item Users add conversation/evidence dynamically (e.g., ``only auto Defendant, manual others'', or ``all auto'').
\item \textbf{Post-trial} (user says ``end''), Case Predictor outputs win probability + feedback/pleadings/replay.
\end{enumerate}

\subsection{Tech Highlights}
\begin{itemize}[leftmargin=*]
\item MARL for lawyer agent realism
\item Thai Civil Code integration
\item Plain-language UI
\end{itemize}


\section{AI and Data Design}
\label{section:ai-data-design}

\subsection{Data Sources and Collection}
<TIP: Describe where your data comes from. Is it collected from users,
scraped from the web, sourced from a public dataset, or
provided by an external API?
Explain how much data you have or plan to collect,
and any challenges you anticipate in obtaining it.
Explain how data is cleaned, transformed, or prepared before use. />
\subsection{Model Design}
% <TIP: Describe the AI model or external AI service you plan to use.
% Explain why you chose it over alternatives.
% This can connect back to your competitor or literature review in Chapter 2.
% If you are training your own model, describe the architecture and
% training approach. If you are using a third-party model or API,
% describe how it is integrated into your system and any prompt design
% or fine-tuning involved. />
Judge:
Lawyer:
Defender:
Case Predictor:

\subsection{Evaluation}
% <TIP: Describe the metrics you plan to use to evaluate your AI model. />
Use case outcome prediction accuracy as the primary metric for evaluating our AI model's performance. This metric measures the percentage of cases where the model's predicted outcome (win, lose, settle) matches the actual outcome of the case. We will also analyze precision and recall for each outcome category to understand the model's strengths and weaknesses in predicting different types of case results.
Also, we'll use the context from the generated text from the fine-tuned lawyer, defender, and judge for the prediction of the case outcome. This will help us evaluate how well the model understands and utilizes the legal context in making its predictions.
